\textbf{Phần 1. Trắc nghiệm lựa chọn}
\setcounterpageref{ex}{0}

\begin{ex} %[5D4N19]
Kết quả của phép tính $6{,}8 + 9{,}7$ là
\choice
{$15{,}5$}
{$16{,}15$}
{$17{,}5$}
{\True $16{,}5$}
\loigiai{
Ta có: $6{,}8 + 9{,}7 = 16{,}5$.\\
\textbf{Đáp án của câu hỏi là D}
}
\end{ex}

\begin{ex} %[5D4N19]
Kết quả của phép tính $24{,}5 + 3{,}84$ là
\choice
{$27{,}34$}
{$28{,}89$}
{\True $28{,}34$}
{$62{,}9$}
\loigiai{
Đặt tính thẳng cột: $24{,}50 + 3{,}84 = 28{,}34$.\\
\textbf{Đáp án của câu hỏi là C}
}
\end{ex}

\begin{ex} %[5D4H19]
Kết quả của phép tính $8{,}16 + 5{,}9$ là
\choice
{\True $14{,}06$}
{$13{,}25$}
{$14{,}6$}
{$13{,}06$}
\loigiai{
Ta có: $8{,}16 + 5{,}90 = 14{,}06$.\\
\textbf{Đáp án của câu hỏi là A}
}
\end{ex}

\begin{ex} %[5D4H19]
Kết quả của phép tính $18{,}6 + 8{,}72$ là
\choice
{$26{,}32$}
{\True $27{,}32$}
{$27{,}78$}
{$105{,}8$}
\loigiai{
Ta có: $18{,}60 + 8{,}72 = 27{,}32$.\\
\textbf{Đáp án của câu hỏi là B}
}
\end{ex}

\begin{ex} %[5D4H19]
Kết quả của phép tính $4{,}61 + 8$ là
\choice
{\True $12{,}61$}
{$5{,}41$}
{$4{,}69$}
{$12{,}69$}
\loigiai{
Cộng phần nguyên: $4 + 8 = 12$, giữ nguyên phần thập phân: $12{,}61$.\\
\textbf{Đáp án của câu hỏi là A}
}
\end{ex}

\begin{ex} %[5D4H19]
Kết quả của phép tính $23{,}58 + 5{,}271$ là
\choice
{$28{,}751$}
{\True $28{,}851$}
{$76{,}29$}
{$28{,}85$}
\loigiai{
Ta có: $23{,}580 + 5{,}271 = 28{,}851$.\\
\textbf{Đáp án của câu hỏi là B}
}
\end{ex}

\begin{ex} %[5D4H19]
Biểu thức $37{,}5 + 4{,}6 = 4{,}6 + 37{,}5$ thể hiện tính chất nào của phép cộng?
\choice
{Tính chất kết hợp}
{Tính chất cộng với số 0}
{\True Tính chất giao hoán}
{Tính chất phân phối}
\loigiai{
Khi đổi chỗ hai số hạng trong một tổng thì tổng không thay đổi (tính chất giao hoán).\\
\textbf{Đáp án của câu hỏi là C}
}
\end{ex}

\begin{ex} %[5D4H19]
Tính thuận tiện biểu thức $6 + 8{,}46 + 1{,}54$ ta được kết quả là
\choice
{$15$}
{$17$}
{$16{,}46$}
{\True $16$}
\loigiai{
Ta có: $6 + (8{,}46 + 1{,}54) = 6 + 10 = 16$.\\
\textbf{Đáp án của câu hỏi là D}
}
\end{ex}

\begin{ex} %[5D4H19]
Tính thuận tiện biểu thức $4{,}8 + 0{,}73 + 3{,}2$ ta được kết quả là
\choice
{\True $8{,}73$}
{$7{,}73$}
{$8{,}37$}
{$9{,}73$}
\loigiai{
Ta có: $(4{,}8 + 3{,}2) + 0{,}73 = 8 + 0{,}73 = 8{,}73$.\\
\textbf{Đáp án của câu hỏi là A}
}
\end{ex}

\begin{ex} %[5D4V19]
Mai cùng mẹ đi siêu thị mua một quả dưa hấu cân nặng $4{,}65\text{ kg}$ và một quả mít cân nặng $5{,}8\text{ kg}$. Cả hai quả dưa hấu và mít cân nặng tất cả là
\choice
{$9{,}45\text{ kg}$}
{$10{,}35\text{ kg}$}
{\True $10{,}45\text{ kg}$}
{$11{,}45\text{ kg}$}
\loigiai{
Cả hai quả cân nặng: $4{,}65 + 5{,}8 = 10{,}45\text{ (kg)}$.\\
\textbf{Đáp án của câu hỏi là C}
}
\end{ex}

\begin{ex} %[5D4V19]
Ngày thứ Sáu chú Sơn thu được $73{,}5\text{ l}$ sữa dê, ngày thứ Bảy thu được nhiều hơn ngày thứ Sáu $4{,}75\text{ l}$ sữa dê. Cả hai ngày chú Sơn thu được số lít sữa dê là
\choice
{$78{,}25\text{ l}$}
{$150{,}75\text{ l}$}
{$152{,}25\text{ l}$}
{\True $151{,}75\text{ l}$}
\loigiai{
Ngày thứ Bảy chú Sơn thu được: $73{,}5 + 4{,}75 = 78{,}25\text{ (l)}$.\\
Cả hai ngày chú Sơn thu được: $73{,}5 + 78{,}25 = 151{,}75\text{ (l)}$.\\
\textbf{Đáp án của câu hỏi là D}
}
\end{ex}

\begin{ex} %[5D4V19]
Sân trường dạng hình chữ nhật có chiều rộng $17{,}5\text{ m}$, chiều dài hơn chiều rộng $15\text{ m}$. Chu vi sân trường đó là
\choice
{$50\text{ m}$}
{\True $100\text{ m}$}
{$65\text{ m}$}
{$130\text{ m}$}
\loigiai{
Chiều dài sân trường là: $17{,}5 + 15 = 32{,}5\text{ (m)}$.\\
Chu vi sân trường là: $(32{,}5 + 17{,}5) \times 2 = 50 \times 2 = 100\text{ (m)}$.\\
\textbf{Đáp án của câu hỏi là B}
}
\end{ex}

\textbf{Phần 2. Trắc nghiệm đúng sai}
\setcounterpageref{ex}{0}

\begin{ex} %[5D4H19]
Xét tính đúng sai của việc thực hiện các phép tính sau:
\choiceTF
{Phép cộng $5{,}4 + 3{,}9 = 8{,}13$.}
{\True Phép cộng $2{,}45 + 4{,}37 = 6{,}82$.}
{Phép cộng $6{,}53 + 12{,}8 = 78{,}1$.}
{\True Phép tính $(5 + 2{,}7) + 7{,}3 = 5 + (2{,}7 + 7{,}3) = 15$.}
\loigiai{
\begin{itemize}
    \item $5{,}4 + 3{,}9 = 9{,}3 \neq 8{,}13$ (sai do cộng từng phần rồi ghép).
    \item $2{,}45 + 4{,}37 = 6{,}82$.
    \item $6{,}53 + 12{,}8 = 19{,}33 \neq 78{,}1$ (sai do đặt tính lệch cột).
    \item $(5 + 2{,}7) + 7{,}3 = 5 + 10 = 15$.
\end{itemize}
\begin{itemize}
    \item a) \textbf{Sai}: $5{,}4 + 3{,}9 = 9{,}3$.
    \item b) \textbf{Đúng}.
    \item c) \textbf{Sai}: $6{,}53 + 12{,}8 = 19{,}33$.
    \item d) \textbf{Đúng}.
\end{itemize}
}
\end{ex}

\begin{ex} %[5D4V19]
Hai bạn Mai và Việt dùng dây đồng để uốn hình. Mai uốn hình ngôi sao hết $1{,}65\text{ m}$ dây đồng, Việt uốn hình tam giác hết $1{,}26\text{ m}$ dây đồng. Xét tính đúng sai của các khẳng định sau:
\choiceTF
{\True Chiều dài dây đồng Mai dùng nhiều hơn của Việt là $0{,}39\text{ m}$.}
{\True Tổng chiều dài dây đồng cả hai bạn đã dùng là $2{,}91\text{ m}$.}
{Nếu có đoạn dây đồng dài $3\text{ m}$, sau khi hai bạn uốn xong thì còn thừa $0{,}19\text{ m}$.}
{\True Đổi $2{,}91\text{ m} = 291\text{ cm}$.}
\loigiai{
\begin{itemize}
    \item Mai dùng nhiều hơn Việt: $1{,}65 - 1{,}26 = 0{,}39\text{ m}$.
    \item Cả hai bạn dùng: $1{,}65 + 1,26 = 2{,}91\text{ m}$.
    \item Đoạn dây còn thừa: $3 - 2{,}91 = 0{,}09\text{ m} \neq 0{,}19\text{ m}$.
    \item $2{,}91\text{ m} = 291\text{ cm}$.
\end{itemize}
\begin{itemize}
    \item a) \textbf{Đúng}.
    \item b) \textbf{Đúng}.
    \item c) \textbf{Sai}: Còn thừa $0{,}09\text{ m}$.
    \item d) \textbf{Đúng}.
\end{itemize}
}
\end{ex}

\textbf{Phần 3. Trả lời ngắn}
\setcounterpageref{ex}{0}

\begin{ex} %[5D4H19]
Tính tổng của hai số $5{,}34$ và $7{,}49$. (Nhập kết quả dưới dạng số thập phân).
\shortans{12,83}
\loigiai{
\textbf{Đáp số:} $12{,}83$\\
$5{,}34 + 7{,}49 = 12{,}83$.
}
\end{ex}

\begin{ex} %[5D4H19]
Tính tổng của hai số $7{,}635$ và $6{,}249$. (Nhập kết quả dưới dạng số thập phân).
\shortans{13,884}
\loigiai{
\textbf{Đáp số:} $13{,}884$\\
$7{,}635 + 6{,}249 = 13{,}884$.
}
\end{ex}

\begin{ex} %[5D4H19]
Tính nhanh giá trị biểu thức: $4{,}8 + 0{,}73 + 3{,}2$. (Nhập kết quả dưới dạng số thập phân).
\shortans{8,73}
\loigiai{
\textbf{Đáp số:} $8{,}73$\\
$(4{,}8 + 3{,}2) + 0{,}73 = 8 + 0{,}73 = 8{,}73$.
}
\end{ex}

\begin{ex} %[5D4V19]
Một quả dưa hấu nặng $4{,}65\text{ kg}$, quả mít nặng $5{,}8\text{ kg}$. Hỏi cả hai quả nặng bao nhiêu ki-lô-gam? (Nhập kết quả dưới dạng số thập phân).
\shortans{10,45}
\loigiai{
\textbf{Đáp số:} $10{,}45$\\
$4{,}65 + 5{,}8 = 10{,}45\text{ (kg)}$.
}
\end{ex}

\begin{ex} %[5D4V19]
Ngày thứ nhất đội công nhân sửa được $73{,}5\text{ m}$ đường, ngày thứ hai sửa được nhiều hơn ngày thứ nhất $4{,}75\text{ m}$ đường. Hỏi cả hai ngày đội sửa được bao nhiêu mét đường? (Nhập kết quả dưới dạng số thập phân).
\shortans{151,75}
\loigiai{
\textbf{Đáp số:} $151{,}75$\\
Ngày hai: $73{,}5 + 4{,}75 = 78{,}25\text{ m}$. Cả hai ngày: $73{,}5 + 78{,}25 = 151{,}75\text{ m}$.
}
\end{ex}

\begin{ex} %[5D4V19]
Một sân trường hình chữ nhật có chiều rộng $17{,}5\text{ m}$, chiều dài hơn chiều rộng $15\text{ m}$. Tính chu vi sân trường đó theo đơn vị mét.
\shortans{100}
\loigiai{
\textbf{Đáp số:} $100$\\
Chiều dài: $17{,}5 + 15 = 32{,}5\text{ m}$. Chu vi: $(32{,}5 + 17{,}5) \times 2 = 100\text{ m}$.
}
\end{ex}
